\section{Response to Reviewer \texttt{QwUG}}

\subsection*{Summary}

The manuscript is well written, clearly structured, and technically rigorous, with detailed experiments and robustness checks. The manuscript would benefit from more discussion of key assumptions in real-world clinical trial settings and clearer description of the experimental data sources.


\subsection*{Summary of Strengths}

\begin{itemize}
    \item The manuscript introduces a placebo-anchored proxy--gold transfer learning framework, which is an innovative way to combine abundant multi-site RCT data (proxy labels) with high-fidelity placebo outcomes from the gold labels.
    \item The manuscript is methodologically rigorous. Stage-wise anchoring using target placebo outcomes, and doubly robust orthogonalization using known randomization propensities, demonstrate careful research design.
    \item The manuscript includes strong experimental design. Well-designed evaluation metrics—Precision in Estimation of Heterogeneous Effects (PEHE), mean absolute error of the average treatment effect (ATE), and calibration RMSE for treated and placebo outcome regressions ($\mu_1$ and $\mu_0$)—capture both individual-level and population-level performance.
    \item Writing is clear and structure is logical.
\end{itemize}

\subsection*{Summary of Weaknesses}

\subsubsection{Justification of Assumptions A4 and A5 in real-world settings.}  
Assumptions A4 and A5, specialized to the multi-site RCT setting, are crucial for performance. However, the manuscript does not sufficiently justify why real-world clinical trials or clinical data would satisfy these assumptions. Additional discussion is suggested.\\

{\color{blue}
We thank the reviewer for this suggestion and agree that additional clinical context would strengthen the manuscript. 
We emphasize that Assumptions~A4 and~A5 are intended as \emph{regularity conditions for stable estimation} in multi-site randomized trials, rather than strong biological invariance or causal identification assumptions.

Assumption~A4 has now been refined into a Lipschitz smoothness condition on outcome regressions: it requires that conditional mean outcomes vary continuously with baseline covariates within each site. 
This assumption rules out arbitrarily unstable responses to small covariate perturbations and ensures that moderate covariate shift across trial populations does not induce unbounded changes in predicted outcomes. 
Such smoothness assumptions are standard in nonparametric regression and causal learning, and are implicitly relied upon whenever outcome models are evaluated outside their exact training distribution.
Importantly, A4 does not assert exchangeability, treatment-effect homogeneity, or transportability across sites.

Assumption~A5 concerns how outcome regressions differ across trial populations. 
Rather than assuming that treatment effects or baseline risks are identical across sites, A5 posits that systematic cross-site differences can be reasonably approximated by a low-complexity correction to a well-learned proxy model. 
Clinically, this reflects the common observation that between-trial heterogeneity is often driven by a limited number of interpretable factors—such as baseline risk calibration, disease severity, key biomarkers, or measurement and enrollment practices—rather than by wholesale changes in the outcome-generating mechanism.

This perspective is consistent with individual participant data meta-analysis and clinical prediction model updating, where cross-study heterogeneity is typically addressed through parsimonious recalibration or limited effect-modification, rather than unrestricted refitting of complex models.
Importantly, Assumption~A5 is \emph{no longer} an identification assumption: it defines a working approximation class that enables stable estimation and calibrated extrapolation when treated outcomes are unavailable in the target population. 
When A5 holds approximately, estimation error decomposes into stochastic error and approximation error; when it is violated, performance degrades smoothly toward proxy-only baselines rather than exhibiting a sharp failure.
}


\subsubsection{Description of experimental data sources.}  
The sources of the experimental data are not clearly described. It is recommended to more explicitly describe the data sources and data processing methods.

{\color{blue}
We thank the reviewer for this suggestion and have substantially expanded the description of our experimental data in the revised manuscript.

\medskip
\noindent\textbf{Data sources.}
The revised manuscript now reports experiments on \emph{two} data sources: (i)~synthetic multi-center RCT data from a controlled DGP (Sections~4 and appendix), and (ii)~real-life covariate data from the Infant Health and Development Program (IHDP) benchmark in a new \textbf{Real-Life Data Experiments} section (Section~5). IHDP provides 25 real baseline covariates with semi-synthetic outcomes and known ground-truth CATEs; we evaluate connected and disconnected regimes over 50 realizations. Below we summarize the synthetic setup; the IHDP setup and results are described in Section~5 of the main manuscript.

\medskip
\noindent\textbf{Synthetic multi-center RCT data.}
Synthetic experiments use data generated from a controlled Data-Generating Process (DGP) designed to mimic realistic multi-site randomized controlled trial settings. We use synthetic data for the following reasons:
\begin{enumerate}[leftmargin=*]
\item \textbf{Ground-truth access}: Synthetic data allows us to compute ground-truth individual treatment effects $\tau(X_i)$, which are unobservable in real data. This enables direct evaluation of CATE estimation accuracy via PEHE and calibration metrics.
\item \textbf{Controlled experimental conditions}: We can systematically vary key parameters (sample size, dimensionality, number of source sites, assumption violations) to understand method behavior across a range of scenarios.
\item \textbf{Reproducibility}: Fixed random seeds ensure fully reproducible experiments.
\end{enumerate}

\medskip
\noindent\textbf{DGP specification (revised Section~IV-C of the manuscript).}
The DGP generates multi-center RCT data with explicit covariate shift, controlled signal-to-noise ratios, and known ground-truth treatment effects:
\begin{itemize}
\item \textbf{Sites}: $C=10$ source sites plus one target site ($c=0$).
\item \textbf{Covariates}: $p$-dimensional baseline covariates $X \in \mathbb{R}^p$ with $p \in \{10, 20, 50, 100\}$. Site-level covariate shift is induced via $X \mid c \sim \mathcal{N}(\mu_c, I)$, calibrated to produce moderate overlap (AUC $\approx 0.75$ for source vs.\ target classification).
\item \textbf{Treatment assignment}: Randomized with logistic propensity $e(X)$, known by design.
\item \textbf{Outcomes}: Follow the proxy--deviation decomposition from Assumption~A5:
$\mu_{a,c}(x) = \mu^{\mathrm{proxy}}_a(x) + \delta_{a,c}(x)$,
where $\mu^{\mathrm{proxy}}_a(x) = x^\top \beta_a$ are shared outcome regressions and $\delta_{a,c}(x) = x^\top \gamma_{a,c}$ are sparse site-specific deviations.
\item \textbf{Sample sizes}: Target budget $(m_0, m_1)$ varies across experiments. Source sites contribute 20,000 total observations (2,000 per site).
\end{itemize}

\medskip
\noindent\textbf{Experimental sweeps.}
We conduct three primary sweeps to evaluate robustness:
\begin{enumerate}[leftmargin=*]
\item \textbf{Target budget vs.\ dimensionality}: $(m_0, m_1) \times p \in \{10, 20, 50, 100\}$
\item \textbf{Source site scaling}: Number of sources $C \in \{2, 5, 10, 20, 50\}$
\item \textbf{Assumption A5 violation}: Sparsity ratio and nonlinearity parameter $\Gamma$ varied to test graceful degradation
\end{enumerate}
All experiments use 100 Monte Carlo replications with fixed seeds, reporting mean $\pm$ standard deviation across replications.

\medskip
\noindent\textbf{Implementation.}
All simulations are implemented in Python. Regularization parameters are selected via cross-validation. The \texttt{glmtrans} transfer learning algorithm is interfaced through an R package wrapper. Code and DGP specifications are available in the supplementary materials.
}


\clearpage